\subsection{Removing the supplied-support condition on arbitrary graphs}
\label{subsec:condition-free-graph}

The response machinery is not required for correctness.  It is required only
to accelerate the following condition-free local gate.  Continue to assume a
single root seed $o$ on an arbitrary graph, and recall
$c_\rho=b-\alpha\rho D^{1/2}\one$.
If $c_{\rho,o}\leq0$, return zero.  Otherwise initialize $U_0=\{o\}$ and, at
stage $j$, compute
\begin{equation}
 Q_{U_jU_j}x^{(j)}_{U_j}=c_{\rho,U_j},
 \qquad x^{(j)}_{U_j^c}=0,
 \label{eq:exact-tree-gate-solve}
\end{equation}
then form the violating boundary batch
\begin{equation}
 B_j
 :=\left\{
  v\in\partial U_j:
  r_v(x^{(j)})>\alpha\rho\sqrt{d_v}
 \right\}.
 \label{eq:exact-tree-gate-batch}
\end{equation}
If $B_j=\emptyset$, return $x^{(j)}$; otherwise set
$U_{j+1}=U_j\cup B_j$.

\begin{theorem}[Condition-free exact local RPPR on an arbitrary graph]
\label{thm:condition-free-graph}
The exact boundary gate is well defined and satisfies, at every nonterminal
stage,
\begin{equation}
 x^{(j)}_{U_j}>0,
 \qquad
 0\leq x^{(j)}\leq x^{(j+1)}\leq x^\star(\rho),
 \qquad
 U_j\subsetneq U_{j+1}\subseteq S^\star(\rho).
 \label{eq:graph-gate-monotonicity}
\end{equation}
It terminates with $U_J=S^\star(\rho)$ and
$x^{(J)}=x^\star(\rho)$ after at most $|S^\star(\rho)|$ nonempty batches.
Writing $J_\star$ for the actual number of nonempty batches, supplied
width-$w$ orderings for the intermediate regions give the sharper
trajectory-dependent bound
\begin{equation}
 \Work_{\rm graph\text{-}gate}
 =O\!\left(
  (w+1)^2(J_\star+1)\vol(S^\star(\rho))
 \right).
 \label{eq:graph-gate-batch-work}
\end{equation}
If every $G[U_j]$ is a forest, fresh elimination and one boundary scan at each
stage give
\begin{equation}
 \Work_{\rm graph\text{-}gate}
 =O\!\left(
  |S^\star(\rho)|\vol(S^\star(\rho))
 \right)
 =O(1/\rho^2).
 \label{eq:graph-gate-work}
\end{equation}
More generally, supplied width-$w$ orderings for the intermediate regions give
$O((w+1)^2/\rho^2)$ work.  No exact support, radius, or confinement oracle is
assumed; the width ordering is needed only for this arithmetic bound, not for
correctness or discovery.
\end{theorem}

\begin{proof}
If $c_{\rho,o}\leq0$, then $x=0$ satisfies every shifted KKT inequality:
the root inequality is the assumed one and every nonseed shifted demand is
strictly negative.  Hence the early return is exact.  When
$c_{\rho,o}>0$, the singleton solution is
$x_o^{(0)}=c_{\rho,o}/Q_{oo}>0$.  Moreover, the active equation of the full
solution gives
\[
 Q_{oo}x_o^\star
 =c_{\rho,o}-Q_{o,S^\star\setminus\{o\}}x_{S^\star\setminus\{o\}}^\star
 \geq c_{\rho,o},
\]
so $x^{(0)}\leq x^\star$.

Suppose inductively that $U_j\subseteq S^\star$,
$x^{(j)}_{U_j}>0$, and $0\leq x^{(j)}\leq x^\star$.  If
$v\notin S^\star$, then $x_v^\star=0$ and the
off-diagonal signs give
\begin{equation}
 \kappa_{\rho,v}(x^{(j)})
 \geq\kappa_{\rho,v}(x^\star)\geq0,
 \qquad
 \kappa_\rho(x):=Qx-b+\alpha\rho D^{1/2}\one.
 \label{eq:tree-safe-admission}
\end{equation}
The inequality in \cref{eq:exact-tree-gate-batch} is equivalent to
$\kappa_{\rho,v}(x^{(j)})<0$.  Hence every admitted vertex belongs to
$S^\star$.

Let $W=U_j\cup B_j$ and
$h_B:=\kappa_{\rho,B_j}(x^{(j)})<0$.  Eliminating the old block gives the
Schur-complement correction
\begin{equation}
 \Delta_B=-\mathcal S_B^{-1}h_B>0,
 \qquad
 \Delta_U=-Q_{UU}^{-1}Q_{UB}\Delta_B\geq0,
 \label{eq:tree-positive-expansion}
\end{equation}
because $\mathcal S_B$ is a principal Schur complement and hence a Stieltjes
matrix with nonnegative inverse.  Thus the enlarged shifted Dirichlet solution
is positive and no smaller than the previous one.  Comparison with the full
solution gives $x^{(j+1)}\leq x^\star$, closing the induction.

Every nonterminal batch is nonempty and lies in the finite set $S^\star$, so
there are at most $|S^\star|$ batches.  At termination, the active equations
hold on $U_J$, the boundary KKT inequalities hold by the empty-batch test, and
every boundary residual is also nonnegative because
$r_v=-Q_{vU_J}x_{U_J}\geq0$.  All more distant nonseed residuals vanish.  The
shifted Dirichlet certificate
\cref{thm:shifted-dirichlet} gives $x^{(J)}=x^\star$, so positivity implies
$U_J=S^\star$.

At stage $j$, \cref{thm:forest-resolution} computes the forest solve and
boundary residual in $O(\vol(U_j))$ work.  With width-$w$ orderings, use
\cref{thm:width-resolution} instead.  Sum over at most $|S^\star|$ stages and
use $U_j\subseteq S^\star$, then
$|S^\star|\leq\vol(S^\star)\leq1/\rho$.
\end{proof}

\begin{theorem}[Proximal domination and margin-sensitive gate bound]
\label{thm:gate-objective-contraction}
On the nonnegative orthant define the shifted obstacle objective and its unit
proximal-gradient map by
\begin{align}
 \Phi_\rho(x)
 &:=\frac12x^\top Qx-c_\rho^\top x
    +\operatorname{ind}_{\mathbb R_+^V}(x),
 &T_\rho(x)&:=[x-Qx+c_\rho]_+.
 \label{eq:shifted-obstacle-objective}
\end{align}
Here $\operatorname{ind}_{\mathbb R_+^V}$ is zero on the nonnegative orthant
and $+\infty$ outside it.
At every nonterminal exact-gate stage,
\begin{equation}
 \Phi_\rho(x^{(j+1)})
 \leq \Phi_\rho(T_\rho(x^{(j)}))
 \leq
 (1-\alpha)\Phi_\rho(x^{(j)})
 +\alpha\Phi_\rho(x^\star).
 \label{eq:gate-objective-contraction}
\end{equation}
Thus, with
$\Delta_j:=\Phi_\rho(x^{(j)})-\Phi_\rho(x^\star)$,
\begin{equation}
 \Delta_j\leq(1-\alpha)^j\Delta_0,
 \qquad
 \Delta_0\leq\frac{\alpha}{2d_o}.
 \label{eq:gate-gap-geometric}
\end{equation}
If $S^\star\neq\{o\}$ and
\begin{equation}
 \delta_\star:=\min_{v\in S^\star}x_v^\star>0,
 \label{eq:active-coordinate-margin}
\end{equation}
then the number of nonempty batches also obeys
\begin{align}
 J_\star
 &\leq
 \min\left\{
 |S^\star|-1,
 1+\left\lceil
 \frac{
  \max\left\{0,\log\!\left(2\Delta_0/(\alpha\delta_\star^2)\right)\right\}}
 {-\log(1-\alpha)}
 \right\rceil
 \right\}
 \notag\\
 &\leq
 \min\left\{
 |S^\star|-1,
 2+\frac{1}{\alpha}
 \max\left\{0,\log\!\frac{1}{d_o\delta_\star^2}\right\}
 \right\}.
 \label{eq:gate-margin-batch-bound}
\end{align}
In particular, the first quantity in \cref{eq:gate-margin-batch-bound} may
replace $J_\star$ in \cref{eq:graph-gate-batch-work}.  This is an
unconditional instance bound: $\delta_\star$ is measured from the unique
optimum, not assumed by the algorithm.
\end{theorem}

\begin{proof}
The active equations give
$T_\rho(x^{(j)})_{U_j}=x^{(j)}_{U_j}$.  Outside $U_j$, a coordinate of the
proximal point is positive exactly when its shifted KKT slack is negative.
A nonboundary nonseed coordinate has no edge to $U_j$ and has strictly
negative shifted demand, so it remains zero.  Consequently
\begin{equation*}
 \supp(T_\rho(x^{(j)}))\subseteq U_j\cup B_j.
\end{equation*}
The next gate point is the exact minimizer of $\Phi_\rho$ on that enlarged
coordinate subspace, proving the first inequality in
\cref{eq:gate-objective-contraction}.

For the second, put $x=x^{(j)}$, $p=T_\rho(x)$, and
$z=(1-\alpha)x+\alpha x^\star$.  The proximal optimality inequality and the
unit smoothness of the quadratic give
\begin{equation*}
 \Phi_\rho(p)
 \leq
 \frac12x^\top Qx-c_\rho^\top x
 +\operatorname{ind}_{\mathbb R_+^V}(z)
 +\langle Qx-c_\rho,z-x\rangle
 +\frac12\|z-x\|_2^2.
\end{equation*}
The orthant indicator is convex.  Applying $\alpha$-strong convexity of the
quadratic between $x$ and $x^\star$ cancels the final
$\alpha^2\|x-x^\star\|_2^2/2$ term and yields the convex combination on the
right side of \cref{eq:gate-objective-contraction}.  Iteration proves the
first part of \cref{eq:gate-gap-geometric}.  Since the singleton gate point
minimizes over a set containing zero, its initial gap is no larger than the
standard single-seed zero-start bound $\alpha/(2d_o)$.

If the gate has not terminated, some $v\in S^\star\setminus U_j$ has
$x_v^{(j)}=0$ and $x_v^\star\geq\delta_\star$.  Strong convexity therefore
gives
\begin{equation*}
 \Delta_j\geq\frac{\alpha}{2}
 \|x^{(j)}-x^\star\|_2^2
 \geq\frac{\alpha\delta_\star^2}{2}.
\end{equation*}
Combining this lower bound with \cref{eq:gate-gap-geometric} proves the first
margin estimate, with one extra batch covering equality.  Finally use
$-\log(1-\alpha)\geq\alpha$ and the initial-gap estimate.  The support-count
bound improves by one because $o\in U_0\subseteq S^\star$.
\end{proof}

\begin{proposition}[A three-path rules out uniform accelerated gap decay]
\label{prop:three-path-gap-obstruction}
No constant $c>0$ can make
\begin{equation}
 \Delta_{j+1}\leq(1-c\sqrt\alpha)\Delta_j
 \label{eq:false-accelerated-gate-gap}
\end{equation}
hold for every exact-gate batch on every single-seed RPPR instance.  Indeed,
on the path $o-u-v$, for all sufficiently small $\rho>0$ the gate visits
$\{o\}$, $\{o,u\}$, and $\{o,u,v\}$, and
\begin{equation}
 \lim_{\rho\downarrow0}\frac{\Delta_1}{\Delta_0}
 =\frac{(1-\alpha)^2}{1+6\alpha+\alpha^2}
 =1-8\alpha+O(\alpha^2).
 \label{eq:three-path-gap-ratio}
\end{equation}
This obstruction concerns a contraction proof, not the batch-count
conjecture: the displayed instance terminates after two nonempty batches.
\end{proposition}

\begin{proof}
At $\rho=0$, write
$a=(1+\alpha)/2$ and $q=(1-\alpha)/(2\sqrt2)$.  In the order $(o,u,v)$,
the matrix has diagonal $a$ and path off-diagonal $-q$, while the demand is
supported at $o$.  A restricted quadratic minimum on the first $k$
coordinates has value
$-\alpha^2(Q_{1:k,1:k}^{-1})_{oo}/2$.  Direct tridiagonal inversion gives
\begin{align*}
 (Q_{1:1,1:1}^{-1})_{oo}&=\frac1a,\\
 (Q_{1:2,1:2}^{-1})_{oo}&=\frac{a}{a^2-q^2},\\
 (Q^{-1})_{oo}&=\frac{a^2-q^2}{a(a^2-2q^2)}.
\end{align*}
Subtracting successive restricted values gives
$\Delta_1/\Delta_0=q^2/(a^2-q^2)$, which simplifies to the rational
expression in \cref{eq:three-path-gap-ratio}.  The strict gate inequalities,
the active-set sequence, and the ratio persist by continuity for all
sufficiently small positive $\rho$.  Since
$1-8\alpha+O(\alpha^2)>1-c\sqrt\alpha$ for sufficiently small $\alpha$,
\cref{eq:false-accelerated-gate-gap} cannot hold uniformly.
\end{proof}

\begin{corollary}[Thin-support product bound]
\label{cor:radial-width-gate}
For a nonzero single-seed problem, define the active radial width
\begin{equation}
 \chi_\star
 :=\max_{0\leq\ell\leq R_\star}
 \left|
  S^\star(\rho)\cap
  \{v:\operatorname{dist}(o,v)=\ell\}
 \right|.
 \label{eq:active-radial-width}
\end{equation}
Under supplied width-$w$ orderings for the intermediate regions, the exact
boundary gate is condition-free and obeys
\begin{align}
 \Work_{\rm radial\text{-}gate}
 &\leq
 O\!\left(
  (w+1)^2\chi_\star(R_\star+1)
  \vol(S^\star(\rho))
 \right)
 \notag\\
 &=\widetilde O\!\left(
  \frac{(w+1)^2\chi_\star}{\rho\sqrt\alpha}
 \right).
 \label{eq:radial-width-gate-work}
\end{align}
Consequently, arbitrarily long cycles and fixed-width cyclic strips meet the
product scale even though their biconnected blocks need not have bounded
size.  On a cycle, \cref{thm:cycle-one-pass} gives the stronger
output-linear bound without a supplied ordering.
\end{corollary}

\begin{proof}
Every nonempty batch admits at least one new true-support vertex, so
\begin{equation*}
 J_\star+1
 \leq |S^\star(\rho)|
 \leq \chi_\star(R_\star+1).
\end{equation*}
Insert this in \cref{eq:graph-gate-batch-work}, then use
$\vol(S^\star)\leq1/\rho$ and the graph-universal radius estimate
\cref{lem:graph-support-radius}.  A cycle has width at most two and at most
two vertices per root-distance shell.  A fixed-width strip has uniformly
bounded elimination width and shell cardinality under its natural layered
ordering.
\end{proof}
